% Part: first-order-logic
% Chapter: axiomatic-deduction
% Section: rules-and-proofs

\documentclass[../../../include/open-logic-section]{subfiles}

\begin{document}

\iftag{FOL}
      {\olfileid[id]{fol}{axd}{rul}}
      {\olfileid[id]{pl}{axd}{rul}}

\olsection{Aturan dan \usetoken{P}{derivation}}

\begin{explain}
  Dalam logika, !!{derivation}s aksiomatik barangkali merupakan sistem
  !!{derivation} yang paling sederhana.
  !!^a{derivation} hanyalah suatu urutan !!{formula}s. Agar terhitung
  sebagai !!a{derivation}, setiap !!{formula} dalam urutan itu harus
  merupakan instansiasi suatu aksioma, atau harus mengikuti satu atau
  lebih !!{formula} yang mendahuluinya dalam urutan melalui suatu aturan
  inferensi. !!^a{derivation} !!{derive} !!{formula} terakhirnya.
\end{explain}

\begin{defn}[!!^{derivability}]
Jika $\Gamma$ adalah himpunan !!{formula}s dari $\Lang L$, maka sebuah
\emph{!!{derivation}} dari~$\Gamma$ adalah urutan berhingga $!A_1$,
\dots,~$!A_n$ dari !!{formula}s sedemikian sehingga, untuk setiap
$i \le n$, salah satu hal berikut berlaku:
\begin{enumerate}
\item $!A_i \in \Gamma$; atau
\item $!A_i$ adalah suatu aksioma; atau
\item $!A_i$ mengikuti dari suatu $!A_j$ (dan $!A_k$) dengan $j < i$
  (dan $k < i$) melalui suatu aturan inferensi.
\end{enumerate}
\end{defn}

Apa yang terhitung sebagai !!{derivation} yang benar bergantung pada
aturan inferensi mana yang kita izinkan (dan tentu saja pada apa yang
kita anggap sebagai aksioma). Suatu aturan inferensi adalah pernyataan
jika-maka yang memberi tahu kita bahwa, dalam kondisi tertentu, sebuah
langkah~$!A_i$ dalam !!a{derivation} merupakan langkah inferensi yang
benar.

\begin{defn}[Aturan inferensi]
Suatu \emph{aturan inferensi} memberikan syarat cukup bagi apa yang
terhitung sebagai langkah inferensi yang benar dalam !!a{derivation}
dari~$\Gamma$.
\end{defn}

Sebagai contoh, karena setiap urutan satu unsur $!A$ dengan
$!A \in \Gamma$ secara langsung terhitung sebagai !!a{derivation},
pernyataan berikut dapat menjadi aturan inferensi yang sangat sederhana:
\begin{quote}
  Jika $!A \in \Gamma$, maka $!A$ selalu merupakan langkah inferensi
  yang benar dalam setiap !!{derivation} dari~$\Gamma$.
\end{quote}
Demikian pula, jika $!A$ merupakan salah satu aksioma, maka $!A$ sendiri
adalah !!a{derivation}, sehingga pernyataan berikut juga merupakan aturan
inferensi:
\begin{quote}
  Jika $!A$ merupakan suatu aksioma, maka $!A$ merupakan langkah
  inferensi yang benar.
\end{quote}
Keadaannya menjadi lebih menarik jika aturan inferensi mengacu pada
!!{formula}s yang muncul sebelum langkah yang sedang diperiksa. Aturan
berikut disebut \emph{modus ponens:}
\begin{quote}
  Jika $!B \lif !A$ dan $!B$ muncul lebih dahulu dalam !!{derivation},
  maka~$!A$ merupakan langkah inferensi yang benar.
\end{quote}
Jika ini merupakan satu-satunya aturan inferensi, maka definisi
!!{derivation} kita di atas berarti sebagai berikut: $!A_1$,
\dots,~$!A_n$ adalah !!a{derivation} jika dan hanya jika, untuk setiap
$i \le n$, salah satu hal berikut berlaku:
\begin{enumerate}
\item $!A_i \in \Gamma$; atau
\item $!A_i$ adalah suatu aksioma; atau
\item untuk suatu $j < i$, $!A_j$ adalah $!B \lif !A_i$, dan untuk
  suatu $k < i$, $!A_k$ adalah~$!B$.
\end{enumerate}
Klausa terakhir mengatakan bahwa $!A_i$ mengikuti dari~$!A_j$
($!B \lif !A_i$) dan $!A_k$ ($!B$) melalui modus ponens. Jika kita dapat
bergerak dari $1$ hingga~$n$, dan setiap kali menemukan
!!a{formula}~$!A_i$ yang berada dalam~$\Gamma$, merupakan suatu aksioma,
atau dinyatakan sebagai langkah inferensi yang benar oleh suatu aturan
inferensi, maka seluruh urutan tersebut terhitung sebagai
!!{derivation} yang benar.

\begin{defn}[!!^{derivability}]
Suatu !!{formula}~$!A$ \emph{!!{derivable}} dari $\Gamma$, ditulis
$\Gamma \Proves !A$, jika terdapat !!a{derivation} dari $\Gamma$ yang
berakhir dengan~$!A$.
\end{defn}

\begin{defn}[Teorema]
!!^a{formula}~$!A$ merupakan suatu \emph{teorema} jika terdapat
!!a{derivation} atas~$!A$ dari himpunan kosong. Kita menulis
$\Proves !A$ jika $!A$ merupakan teorema dan $\Proves/ !A$ jika bukan.
\end{defn}


\end{document}
